\part{Distributed Oscillatory Memory}

\chapter{Empirical Foundations of Macroscale Resonance}\label{ch:cabral}

The empirical motivation for treating distributed temporal coordination as a candidate substrate of memory comes from the study of Cabral, Fernandes, and Shemesh, which recorded functional magnetic resonance imaging from rat brains at a temporal resolution far finer than is customary and analyzed the resulting signals as superpositions of oscillatory modes \cite{cabral2023}. The design, as reported in the published article, was as follows. Six genetically similar female rats were each scanned twice in each of three conditions, giving twelve scans per condition, with ten-minute scans acquired at 9.4 tesla with a temporal resolution of 38 milliseconds, a sampling rate of $26.3$ Hz. The three conditions were defined by isoflurane concentration on a background of medetomidine sedation: none, which the authors call sedation, one percent, called light anesthesia, and three percent, called deep anesthesia. Signals from about fourteen hundred brain voxels in a coronal slice were band-pass filtered below $0.3$ Hz, and scans of a post-mortem rat served as a baseline against which power and component variance were judged. Principal component analysis of the covariance of the filtered signals yielded, under sedation, ten components above the post-mortem baseline whose spatial patterns showed fixed phase relationships across distinct cortical and subcortical subsystems, varying gradually and symmetrically across space, with regions oscillating in phase or in antiphase with each other, which the authors describe as characteristic of standing waves. The temporal signatures of these components, the time courses of their amplitudes, showed transient oscillations with fluctuating amplitude, and the recorded signals could be reconstructed as the linear superposition of the ten components scaled by their temporal signatures. Power extended up to about $0.3$ Hz, peaked between $0.20$ and $0.25$ Hz in cortical voxels under medetomidine and between $0.10$ and $0.15$ Hz in the striatum, and was modulated by anesthetic level. With one percent isoflurane only six components remained above baseline, with altered spatial configuration, shorter-lived oscillations, and power peaking at lower frequencies, and at three percent the variance above baseline was captured by a single component in which cortex and striatum moved together in phase at very low frequency, consistent with the authors' statement that only global aperiodic fluctuations remain under deep anesthesia. The resonance quality factor of the oscillations, which counts the cycles before the autocorrelation decays to $1/e$, was highest under medetomidine and decreased with isoflurane. An extension to a whole-brain acquisition with lower temporal resolution, repetition time $350$ ms across twelve slices in three rats, observed oscillations with overlapping phase relationships.

The authors interpret the findings with a mechanistic model in which the spatial modes are damped resonators operating in the subcritical range of a Hopf bifurcation and driven by background noise, the stochastic resonance of standing waves, whose superposition reproduces the sliding-window correlation structure of long-range functional connectivity. This regime is relevant to the specification below, because it is a damped, noise-driven regime in which the uniform state is stable and oscillations are transient responses, and not the self-sustained regime in which a continuum oscillator model would oscillate autonomously. The distinction is taken up in Chapter~\ref{ch:cgl}.

The observations are strong evidence for one claim and no evidence for several that are adjacent to it. They support the statement that spontaneous activity measured by fMRI in sedated rats is well described by a small repertoire of spatial patterns whose amplitudes oscillate transiently, with phase relationships that are stable within a condition and that depend on the sedative state. They do not establish that those patterns carry information from one region to another, which is a claim about causal transmission that requires interventions or directed measures the study does not provide. They do not establish a binding mechanism, which would require showing that the phase relationships select and group the content represented in the participating regions. They do not establish a role in memory: the recordings were made at rest, under sedation or anesthesia, in six animals, and no encoding or recall of any episode was involved. The authors themselves note that the study did not attempt to separate neural activity from neurovascular coupling. They observe that the temporal signatures do not follow the canonical hemodynamic response, that the frequencies lie above the range typically associated with the blood-oxygenation-level-dependent signal and below cardiac and respiratory rhythms, and that earlier work found coherence between local field potentials and fMRI in this range, and they conclude that a more direct neuronal contribution cannot be excluded. The relationship of these rhythms of a few tenths of a hertz to the theta and gamma rhythms usually invoked in models of hippocampal and cortical memory \cite{buzsaki2006,lisman2013}, which are more than an order of magnitude and about two orders of magnitude faster respectively, is therefore an open empirical question and not an identification. The components obtained by principal component analysis are, finally, defined by variance and orthogonality, which are properties of the method, so the interpretation of each component as a physical standing wave needs the additional tests discussed in Chapter~\ref{ch:spectral}, some of which the authors carry out through the post-mortem control and the comparison with spatially randomized patterns.

For the purposes of this monograph the study is treated as motivation of a modeling stance: macroscale activity can be organized by a small number of coordinated oscillating patterns whose phase relationships are stable and condition-dependent, and the number, frequency, and damping of the modes are state-dependent quantities. The stance suggests that a memory architecture should be allowed to use phase coordination across distributed components, and it identifies the quantities a model should expose to experiment, namely mode counts, mode frequencies, quality factors, and phase relationships, as functions of a state variable analogous to anesthetic depth. It does not constrain the mechanism by which such coordination would subserve recollection, and every architectural claim that goes beyond the stance is labeled in Chapter~\ref{ch:explains} as a hypothesis.

\chapter{Spectral Decomposition of Distributed Activity}\label{ch:spectral}

Let a distributed field be given on a finite index set $\mathcal X$ of locations by
\begin{equation}\label{eq:modes}
\Psi(x,t)=\sum_{k=1}^{K}a_k(t)\,\varphi_k(x),
\end{equation}
where the $\varphi_k$ are spatial modes and the $a_k$ are time-dependent coefficients. Which modes to use and how to estimate them are separate questions with separate answers. Laplacian eigenmodes are obtained from a graph or mesh on $\mathcal X$ with symmetric nonnegative weights $W$, degree matrix $D$, and Laplacian $L=D-W$. The matrix $L$ is symmetric positive semidefinite and has an orthonormal eigenbasis $L\varphi_k=\lambda_k\varphi_k$ with $0=\lambda_1\le\lambda_2\le\cdots$, and the graph Fourier transform of a signal $f$ on $\mathcal X$ is the vector of coefficients $\hat f_k=\ip{\varphi_k}{f}$. A wave equation $\partial_t^2\Psi=-c^2L\Psi$ on the graph has the standing-wave solutions $\Psi(x,t)=\varphi_k(x)\cos(c\sqrt{\lambda_k}\,t+\alpha)$, each with the fixed spatial profile of an eigenmode and angular frequency $c\sqrt{\lambda_k}$, so geometry alone supplies a family of candidate resonant modes whose frequencies are determined by the spectrum of the connectivity. Principal component analysis, in contrast, chooses spatial modes that maximize explained variance of the observed data, with no reference to a generative dynamics. Dynamic mode decomposition fits a linear propagator $x_{t+1}\approx Ax_t$ to the data and uses the eigenpairs $(\mu_j,\phi_j)$ of $A$, with $\mu_j=e^{(\sigma_j+\iu\omega_j)\Delta t}$, so that the growth rate $\sigma_j$ and the angular frequency $\omega_j$ of each mode are read directly from the eigenvalue and the spatial structure, including its phase map, from the eigenvector.

The distinction that matters for the use of the word oscillatory is the one between standing and traveling spatial organization.

\begin{proposition}[Standing versus traveling modes]\label{prop:standing}
Let $\phi=u+\iu v\in\C^N$ and $x(t)=\Re(\phi\,e^{\iu\omega t})=u\cos\omega t-v\sin\omega t$. Then $x(t)$ is a standing wave, meaning that $x(t)$ lies on a fixed line through the origin for all $t$, if and only if $\phi=e^{\iu\alpha}w$ for a real vector $w$ and a real angle $\alpha$, equivalently if and only if the phase differences $\arg\phi_j-\arg\phi_k$ lie in $\{0,\pi\}$ modulo $2\pi$ for all $j,k$ at which both are nonzero. If $u$ and $v$ are linearly independent, $x(t)$ traverses an ellipse in the plane they span and the spatial pattern changes shape over the period.
\end{proposition}

\begin{proof}
If $\phi=e^{\iu\alpha}w$ then $x(t)=w\cos(\omega t+\alpha)$, which lies on the line through $w$. Conversely, if $x(t)$ stays on a line for all $t$ then $u$ and $v$ are linearly dependent, so $v=\kappa u$ or $u=0$, and then $\phi=u(1+\iu\kappa)=\sqrt{1+\kappa^2}\,e^{\iu\arctan\kappa}u$, or $\phi=\iu v$, each of the stated form. The phase-difference statement holds because $\phi_j/\phi_k=w_j/w_k$ is real. When $u,v$ are independent, $x(t)$ is a nondegenerate ellipse.
\end{proof}

A fixed phase relationship across distant locations that is either zero or $\pi$, as reported for the modes of the rat study, is therefore exactly the signature of a standing wave in this sense, and a traveling wave requires a phase map that varies continuously across locations. A principal component pair in quadrature is the typical PCA signature of a traveling wave, since a single traveling wave occupies two components with equal variance and a quarter-period offset.

The legitimacy of calling a spatial mode an oscillatory mode requires conditions beyond the existence of the decomposition. Principal components of any stationary random process with a spectrum concentrated near a frequency will show peaked power spectra, so a spectral peak in a component is not by itself evidence of a generating oscillator. The conditions adopted here are four. The mode must correspond to a complex-conjugate pair of eigenvalues of a fitted propagator with $|\mu|$ close to one, and with $\omega$ separated from its neighbors in the spectrum by more than the estimation error. The mode must be reproducible across disjoint data segments within a condition. The spectral peak must exceed that of surrogate data that preserve the second-order statistics of the original, for example the covariance matrix and the power spectrum of each channel, while randomizing the cross-channel phase relationships. And the spatial phase map must conform to the standing-wave criterion of Proposition~\ref{prop:standing} within estimation error, or be explicitly reported as traveling. Only modes passing these tests are used in Mem8 as oscillatory modes, and modes failing them are used, if at all, as statistical components.

\chapter{Coupled Oscillators and Phase Synchronization}\label{ch:kuramoto}

A minimal model of distributed phase coordination assigns a phase $\theta_i$ to each of $N$ oscillators and couples them by a sinusoidal interaction,
\begin{equation}\label{eq:kuramoto}
\dot\theta_i=\omega_i+\sum_jK_{ij}\sin(\theta_j-\theta_i-\delta_{ij})+u_i(t),
\end{equation}
where $\omega_i$ is a natural frequency, $K_{ij}$ a coupling strength, $\delta_{ij}$ a phase lag, and $u_i$ an external drive. For the case of identical natural frequencies, symmetric coupling $K_{ij}=K_{ji}$, no lag, and no drive, the model is a gradient system, and its behavior is governed by a potential.

\begin{theorem}[Gradient structure of the Kuramoto system]\label{thm:kuramoto-grad}
Let $\omega_i\equiv\omega$, $\delta_{ij}=0$, $u\equiv0$, and $K$ symmetric. In the rotating frame $\vartheta_i=\theta_i-\omega t$, system \eqref{eq:kuramoto} is $\dot\vartheta=-\nabla V(\vartheta)$ with
\begin{equation}\label{eq:V}
V(\vartheta)=-\frac12\sum_{i,j}K_{ij}\cos(\vartheta_j-\vartheta_i).
\end{equation}
Consequently $V$ is nonincreasing along trajectories, with $\dot V=-\sum_i\dot\vartheta_i^2$, every trajectory converges to the set of critical points of $V$, and, since $V$ is real-analytic, to a single critical point. For the explicit Euler discretization $\vartheta^{k+1}=\vartheta^k-\Delta t\,\nabla V(\vartheta^k)$, $V$ is nonincreasing whenever $\Delta t<1/\max_i\sum_j|K_{ij}|$.
\end{theorem}

\begin{proof}
Differentiating, $\partial V/\partial\vartheta_i=-\sum_jK_{ij}\sin(\vartheta_j-\vartheta_i)$, using the symmetry of $K$ to combine the two terms in which $\vartheta_i$ appears. Hence $\dot\vartheta_i=-\partial V/\partial\vartheta_i$ and $\dot V=\sum_i(\partial V/\partial\vartheta_i)\dot\vartheta_i=-\sum_i\dot\vartheta_i^2\le0$. The state space is a torus, hence compact, so the omega-limit set is nonempty and consists of critical points by LaSalle's principle, and convergence to a single point follows from the \L ojasiewicz gradient inequality for analytic potentials. For the discretization, the Hessian of $V$ has entries bounded by Gershgorin's theorem, giving $\norm{\nabla^2V}\le2\max_i\sum_j|K_{ij}|=:\Lambda$, so $\nabla V$ is $\Lambda$-Lipschitz, and the descent lemma gives $V(\vartheta^{k+1})\le V(\vartheta^k)-\Delta t(1-\Delta t\Lambda/2)\norm{\nabla V(\vartheta^k)}^2$, which is a decrease for $\Delta t<2/\Lambda=1/\max_i\sum_j|K_{ij}|$.
\end{proof}

The theorem is the structural basis for the phase block of the minimal model in Appendix~\ref{app:minimal}, whose update is exactly the Euler step above with zero natural-frequency spread. A deterministic check on a random symmetric coupling matrix with $N=30$ oscillators and a step $0.9$ times the bound found $V$ nonincreasing at every one of $400$ steps and the order parameter equal to one to numerical precision at the end (Appendix~\ref{app:minimal}); this checks the statement and does not test any memory property.

Synchronization is not automatically beneficial. The all-to-all case with identical frequencies and uniform positive coupling $K_{ij}=K/N$ is the extreme: it is known that the fully synchronized state is the only stable equilibrium and that the set of initial conditions not converging to it has measure zero \cite{kuramoto1984,strogatz2000,acebron2005,watanabe1994}. In that regime the final phase pattern is the same for every generic initial pattern, so the map from initial relative phases to final relative phases carries no information about the initial pattern, and any code that stores differences among traces in the relative phases of the oscillators is erased. Excess global coupling therefore reduces selectivity, because it collapses the very distinctions the phases were meant to carry, and it increases interference, because a single global phase pattern is shared by every trace read through it. The opposite extreme, no coupling, leaves the phases independent and provides no coordination. The useful regime lies between: structured coupling $K_{ij}$ that is strong within coherence clusters and weak between them, so that the dynamics has several stable or metastable cluster states, each of which can serve as a context for a different subset of traces. Metastability, meaning the persistence of one cluster configuration for a long time followed by a transition to another, arises in the presence of heterogeneity of natural frequencies, delays, or noise, and the transitions between cluster states are the mechanism by which the system can move among dynamically accessible memory states. The monograph measures cluster structure by order parameters of subsets, $R_{\mathcal S}=|\frac1{|\mathcal S|}\sum_{i\in\mathcal S}e^{\iu\theta_i}|$, and by pairwise phase-locking values between oscillators.

\chapter{Oscillatory Modes as a Computational Substrate}\label{ch:cgl}

A continuum model of coordinated oscillation is supplied by the complex Ginzburg--Landau equation, written here with a drive $I$ as
\begin{equation}\label{eq:cgl}
\partial_t\Psi=(D+\iu\Omega)\Delta\Psi+(\alpha+\iu\omega)\Psi-(\beta+\iu\chi)|\Psi|^2\Psi+I(x,t).
\end{equation}
The equation is treated as a proposed Mem8 model and not as an equation established by the neuroscience; its coefficients are free parameters of the specification, and the study discussed in Chapter~\ref{ch:cabral} gives no basis for assigning them. The equation is the normal form near a Hopf bifurcation of a spatially extended medium \cite{aranson2002}, which is the reason it is a natural candidate substrate: it is the generic description of a distributed system of weakly nonlinear oscillators with diffusive coupling.

Two regimes of \eqref{eq:cgl} need to be kept apart because the empirical study fixes only one of them. For $\alpha>0$ the uniform state is unstable and the medium oscillates autonomously, which is the supercritical regime treated in Proposition~\ref{prop:planewave}. For $\alpha<0$ the uniform state is stable, the linearization about it has damped oscillatory modes with decay rates set by $\alpha-Dk^2$ and frequencies set by $\omega-\Omega k^2$, and oscillations appear only as transient responses to a drive $I$ or to noise, with a quality factor that grows as $\alpha\to0^-$. The model of Cabral and colleagues is in the second regime, a subcritical range of a Hopf bifurcation with noise-driven resonance \cite{cabral2023}, so the supercritical regime is a design choice of Mem8 for which the study offers no support, while the damped regime is the one in which a cue can be treated as a resonant drive and the reconstruction as the ringing of the mode that the cue excites. Both are carried through the specification, and the experiments of Chapter~\ref{ch:ablation} include the damped regime as a baseline.

Four features of its solution set bear on memory. The first is the family of plane-wave solutions.

\begin{proposition}[Plane waves]\label{prop:planewave}
For $I\equiv0$, $D>0$, and $\beta>0$, equation \eqref{eq:cgl} has the solutions $\Psi(x,t)=A\,e^{\iu(k\cdot x-\nu t)}$ with
\begin{equation}
A^2=\frac{\alpha-Dk^2}{\beta},\qquad \nu=\Omega k^2-\omega+\chi A^2,
\end{equation}
which exist when $Dk^2<\alpha$. The case $k=0$ is a spatially uniform oscillation with $A^2=\alpha/\beta$.
\end{proposition}

\begin{proof}
For this ansatz $\Delta\Psi=-k^2\Psi$ and $|\Psi|^2=A^2$. Substituting, the left side is $-\iu\nu\Psi$ and the right side is $[-(D+\iu\Omega)k^2+(\alpha+\iu\omega)-(\beta+\iu\chi)A^2]\Psi$. The real part gives $0=-Dk^2+\alpha-\beta A^2$, and the imaginary part gives $-\nu=-\Omega k^2+\omega-\chi A^2$.
\end{proof}

The second feature is stability. After rescaling, the uniform oscillation is stable to long-wavelength perturbations when $1+(\Omega/D)(\chi/\beta)>0$, the Benjamin--Feir--Newell criterion \cite{aranson2002}; when the criterion fails the uniform state gives way to phase turbulence or to defect-mediated turbulence, which would correspond to unstable reconstruction. The third is the structure of the non-oscillatory subcase. When $\Omega=\chi=\omega=0$ and $I$ is time-independent, equation \eqref{eq:cgl} is the gradient flow $\partial_t\Psi=-\delta\mathcal F/\delta\overline\Psi$ of the energy
\begin{equation}\label{eq:GLenergy}
\mathcal F[\Psi]=\int\Bigl(D|\nabla\Psi|^2-\alpha|\Psi|^2+\tfrac\beta2|\Psi|^4-\bigl(\overline I\Psi+I\overline\Psi\bigr)\Bigr)dx,
\end{equation}
whose local minima are stationary attractors, so the non-oscillatory medium behaves as an energy-minimizing associative memory of Hopfield type, with fixed points playing the role of stored patterns and the drive $I$ playing the role of a cue. The oscillatory terms $\Omega,\chi,\omega$ are the non-gradient terms: they are what allows propagating waves, limit cycles, and phase-coded states, and they are also what removes the guarantee that the dynamics settles. Whether these terms improve reconstruction under matched resources is exactly the question posed by the ablation experiments, and the formal separation just given makes the experiment well posed, since the nonoscillatory implementation required in Chapter~\ref{ch:oscengine} is the gradient subcase. The fourth feature concerns localized solutions and boundary conditions: with a spatially localized drive the medium supports localized structures, which would represent individual traces, and the choice of boundary condition, periodic, no-flux, or absorbing, determines which modes are admissible in the way that the geometry determines Laplacian eigenmodes in Chapter~\ref{ch:spectral}. Mode competition arises when two or more modes have positive growth rates and nonlinear saturation couples them, so that the amplitude of one suppresses the others; this is the continuum analogue of the softmax selection used in the discrete architecture.

\chapter{Temporal Binding and Phase-Coded Associations}\label{ch:binding}

A memory of an episode requires that items encountered at different times be associated in a way that preserves what is needed about their temporal relationship. Four relationships need to be distinguished because they are routinely run together. Temporal order is the relation $t_a<t_b$. Simultaneity is $|t_a-t_b|<\epsilon$ for a resolution $\epsilon$. Causal dependence is a relation of intervention, in which the occurrence of $a$ changes the occurrence or character of $b$. Mere coactivation is that $a$ and $b$ were represented at the same time in the system, for any reason. The first two are functions of timestamps, whereas the third and fourth are not, and the difference has a precise consequence: no encoding that depends only on timestamps can distinguish two histories that differ only in their causal structure. Two histories in which $a$ causes $b$ and in which a common cause produces both at the same times have identical timestamps, so by the sufficiency criterion of Chapter~\ref{ch:ident} a timing-only memory cannot support the recollection function that returns the causal structure.

Phase offers a way to encode temporal order and interval. Associate with each item a phase variable that advances at a carrier angular frequency $\omega$, and record at encoding the phase of the item. The phase difference between two items encoded at times $t_a$ and $t_b$ is $\omega(t_b-t_a)$ modulo $2\pi$, so the lag is recovered only modulo the period $T=2\pi/\omega$: a single frequency yields an aliased measurement of lag. The aliasing is resolved by a bank of frequencies.

\begin{proposition}[Multifrequency phase code]\label{prop:crt}
Let lags be integers $\ell\in\{0,\dots,L-1\}$ in units of a time step, and let $T_1,\dots,T_K$ be pairwise coprime integer periods with $\prod_kT_k\ge L$. The phase differences $\omega_k\ell\bmod 2\pi$ with $\omega_k=2\pi/T_k$ determine $\ell$ uniquely. Conversely, for any set of periods whose least common multiple is smaller than $L$, two distinct lags yield identical phase codes.
\end{proposition}

\begin{proof}
The phase difference at frequency $\omega_k$ determines $\ell\bmod T_k$. By the Chinese remainder theorem the residues $\ell\bmod T_k$ for pairwise coprime moduli determine $\ell$ modulo $\prod_kT_k$, hence uniquely in $\{0,\dots,L-1\}$. If $\mathrm{lcm}(T_1,\dots,T_K)=\Lambda<L$, then lags $0$ and $\Lambda$ both lie in $\{0,\dots,L-1\}$ and agree modulo every $T_k$.
\end{proof}

Static positional encodings in sequence models, which use sines and cosines at a geometric ladder of frequencies \cite{vaswani2017}, are phase codes of exactly this kind, assigned by the position index and not generated by dynamics. This observation fixes the right test for whether phase carries information beyond a static positional encoding. A memory in which the phases are produced by interaction of the oscillators with the input history, and in which retrieval is read from the relative phases, must be compared with a memory that receives the same positional codes as inputs. Two controls are specified. A phase-shuffle control randomizes the assignment of phases to events while preserving their marginal distribution; if sequence reconstruction is unchanged, phase was not being used. A frozen-phase control fixes the oscillator phases to their initial values; if performance is unchanged, the dynamics were not contributing. These are the controls on which the ablation program of Chapter~\ref{ch:ablation} is based.

The reconstruction of an ordered sequence from such a code proceeds by reading relative phases of the stored items at the carrier frequencies and sorting by the decoded lag. The decoded lag error under phase noise of standard deviation $\sigma_k$ at frequency $\omega_k$ is of order $\sigma_k/\omega_k$ for the lowest frequency that resolves the coarse order, and the highest frequency in the bank sets the fine resolution but is the most sensitive to noise in the same ratio, which gives the standard coarse-to-fine tradeoff of residue-number codes. The experimental protocol of Chapter~\ref{ch:sequence} reports reconstruction error as a function of phase noise for this reason.
